% ---
% id: meat-core/def:classes
% kind: def
% title: "Sector classes and the removal checklist"
% summary: >-
%   Exact predicates on the declared data make a sector key-active,
%   constraint-active, state-blind or removable; for a list without
%   party-local entries, removability is certified by a witness meeting a
%   six-condition checklist. Deleting several removable sectors needs a
%   further condition on the set.
% prominence: standard
% status: implemented
% uses: [def:keyinert, prop:exact-cancellation, def:feasible]
% mentions: [thm:removal, prop:removal-compose, def:declaration, def:joint-sectors, rem:target-conditions, rem:witness-exactness, rem:removal-sets, foundations/def:binary64, def:admissible-declaration]
% assumes: []
% refs: []
% provenance: newly-derived
% implements: certiqs_key.meat_core.declaration.sectors:classify_sector
% lean:
% notes: "The item letters (a)-(d) and the checklist labels (W1)-(W6) are cited by letter from other claims (prop:removal-compose, rem:removal-sets, the removal requirements); keep them."
% ---

Two operations on the declared sectors cost nothing: evaluating the
objective over the sectors that carry a deficit, and deleting a sector whose
weight other key-inert sectors can take over. The classes below name the
sectors each operation applies to, and the checklist says when a sector may
be deleted. Neither half of (W2) is redundant
(Remark~\ref{rem:target-conditions}); why (W3) and (W4) must hold exactly,
and why (W4) names one form, is Remark~\ref{rem:witness-exactness}; and
Remark~\ref{rem:removal-sets} exhibits two removable sectors whose joint
deletion overstates the value.

\begin{definition}\label{def:classes}
Let $\Decl$ satisfy Definition~\ref{def:admissible-declaration}. For $s\in\SecSet$:
\begin{enumerate}[label=(\alph*),leftmargin=2.2em]
\item $s$ is \term[Said of a sector that fails to be key-inert for some key-live acceptance value, so that it carries entropy deficits in the objective.]{key-active} if $(\hat c,s)$ fails to be key-inert
(Definition~\ref{def:keyinert}) for at least one $\hat c\in\Chkey$; these
sectors carry the deficits $f_{\hat c,s}$ in the objective.
\item $s$ is \term[Said of a sector that is key-inert for every key-live acceptance value: it adds no deficit but remains a variable of the program.]{constraint-active} if it is key-inert for every
$\hat c\in\Chkey$; by Proposition~\ref{prop:exact-cancellation} it contributes nothing to
any $K_{\hat c}$ but remains an optimization variable (it enters
$\Phimap$, every $P_{\hat c}$, the marginal constraint, and $\CList$).
\item $s$ is \term[Said of a sector on which every measurement operator of the measured system, and every constraint operator when constraints are declared, is a scalar multiple of the identity.]{state-blind} if $F^{(s)}_v=c_{v,s}\,\idon{\Qsys_s}$ for
scalars $c_{v,s}\ge0$, for every $v\in\Vset$ (then
$\sum_vc_{v,s}=1$ by POVM completeness restricted to the sector); if
$\CList\neq\varnothing$ we additionally require each
$\hat O_k^{(s)}=o_{k,s}\,\idon{\Qsys_s}$ for scalars $o_{k,s}$.
\item $s$ is \term[Said of a sector for which the hypotheses of the sector-removal theorem hold, as certified by a witness meeting the removal checklist.]{removable} if the hypotheses of
Theorem~\ref{thm:removal} hold for it.
\end{enumerate}
All four predicates are decided by finitely many linear-algebra checks on
the declaration. Items (a) and (b) reduce to key-inertness
(Definition~\ref{def:keyinert}), a finite combinatorial predicate on the
declared function tables and on which of the $M_x$, $F^{(s)}_v$ are the zero
operator, decided exactly, at no window. Items (c) and (d) are decided
exactly as well: (c) is a condition on the
\tref[declared value]{declared values} themselves
(\xref{foundations/def:binary64}), so that a block that is only
\emph{nearly} scalar fails it and $s$ is then not state-blind, and (d) is
decided by the checklist below.

\emph{The removal checklist.} A \emph{candidate witness} for a sector
$s^{*}$ is a target list $s_1,\dots,s_L$ as index data, states $\chi_i$ as
binary64 sector-block data, weights $\mu_i$ as binary64 data, and the
evaluation order of (W6). It is a \emph{witness} for $s^{*}$ in the sense
of~(d) when $\CList$ carries no entry declared \tref{party-local} and all six
of the following hold:
\begin{enumerate}[label=(W\arabic*),leftmargin=3.0em]
\item \emph{$s^{*}$ is state-blind} in the sense of (c), the
$\hat O^{(s^{*})}_k$ half included whenever $\CList\neq\varnothing$;
\item \emph{the witness does not name $s^{*}$, and names no sector
twice}: $L\ge1$, $s_i\in\SecSet\setminus\{s^{*}\}$ for every $i$, and the
$s_i$ pairwise distinct, all compared as index data at
no window. A candidate naming $s^{*}$ fails (W2), and so does one that
repeats a target;
\item \emph{weights}: $\mu_i\ge0$ for every $i$ --- a comparison of an
exhibited binary64 datum against zero --- and $\sum_i\mu_i=1$, the sum
formed in the evaluation order the witness exhibits and verified to
\tref[rounds nowhere]{round nowhere}, exactly as (W6) requires of its
residuals, so that the computed sum \emph{is} the real sum;
\item \emph{states}: $\chi_i\succeq0$ and $\tr\chi_i=1$ for every $i$,
decided on the \emph{exhibited form} and never on a computed spectrum:
the witness presents each $\chi_i$ diagonal in the declared basis of
$\Qsys_{s_i}$ --- every off-diagonal entry the datum $0$ in the sense of
(c), every diagonal entry a datum $\ge0$, and their sum formed in the
exhibited order, rounding nowhere, and equal to the datum $1$ --- so that
both conditions hold of the declared numbers by construction. Any other
presentation fails (W4), which can only \emph{refuse} a removal;
\item \emph{key-inert targets}, which is (RM2): every pair
$(\hat c,s_i)$ with $\hat c\in\Chkey$ and $i\in\{1,\dots,L\}$ is key-inert
(Definition~\ref{def:keyinert});
\item \emph{both equality families close at residual exactly zero}: (RM1)
for every $v\in\Vset$, and (RM3) for every $k$ (vacuous when
$\CList=\varnothing$).
\end{enumerate}
The condition on $\CList$ and (W1)--(W6) are the hypotheses of
Theorem~\ref{thm:removal} transcribed, all of them and in the theorem's own
words --- the condition on $\CList$ from its first hypothesis, (W1) and
(W2) from its opening sentence, (W3) and (W4) from its quantifiers over
$\mu_i$ and $\chi_i$, (W5) from (RM2), (W6) from (RM1) and (RM3) --- which
is what (d) means by ``the hypotheses of Theorem~\ref{thm:removal}
hold''.
Condition (W6) is read on the residual itself: each residual is formed in
an evaluation order that the witness exhibits alongside its data, every
elementary operation of which rounds nothing, so that the computed residual
\emph{is} the real residual and its comparison against zero decides the
real identity.

\emph{Deleting a set of removable sectors needs one condition more, and it is
not part of this definition.} Removability is a property of one sector, and
Theorem~\ref{thm:removal} licenses the deletion of one sector; that two
sectors are each removable does not license deleting both.
Proposition~\ref{prop:removal-compose} states the further hypothesis on the
set $\SecSet_{\mathrm{del}}$ a run deletes.
Nothing in (W1)--(W6) implies that hypothesis.

\emph{In the general format} of Definition~\ref{def:declaration} the four
classes remain predicates of the sectors $s\in\SecSet$ of $\Qsys$, decided
as above, with $\hat O^{(s)}_k$ the block of the \tref{lifted operator} of an
entry. The finer notion along Alice's sectors --- which joint
pairs $(\hat c,(t,s))$ carry a deficit --- is
Definition~\ref{def:joint-sectors}, and a sector that is key-active here may
have key-inert joint pairs. Theorem~\ref{thm:removal} is stated for a list
carrying no entry declared party-local, so for a list carrying one,
clause~(d) is unsatisfiable and no sector carries a witness in its sense.
\end{definition}
